\documentclass[11pt,a4paper]{article}

\usepackage[utf8]{inputenc}
\usepackage[T1]{fontenc}
\usepackage{lmodern}
\usepackage{microtype}
\usepackage[margin=2.5cm]{geometry}
\usepackage{amsmath,amssymb,amsthm}
\usepackage{mathtools}
\usepackage{booktabs}
\usepackage{longtable}
\usepackage{array}
\usepackage{enumitem}
\usepackage{hyperref}
\usepackage{xcolor}
\usepackage{tcolorbox}
\usepackage{fancyhdr}
\usepackage{titlesec}
\usepackage{csquotes}
\usepackage{listings}

\lstset{
  basicstyle=\ttfamily\small,
  breaklines=true,
  columns=fullflexible,
  showstringspaces=false,
  frame=single,
  framerule=0.2pt,
  rulecolor=\color{gray!30},
  backgroundcolor=\color{gray!2}
}

\hypersetup{
  colorlinks=true,
  linkcolor=blue!50!black,
  citecolor=blue!50!black,
  urlcolor=blue!50!black,
}

\pagestyle{fancy}
\fancyhf{}
\fancyhead[L]{\small The Translation Layer of the V$_{600}$ Substrate}
\fancyhead[R]{\small Pre-peer-review preprint}
\fancyfoot[C]{\thepage}

\newtheorem{definition}{Definition}
\newtheorem{theorem}{Theorem}
\newtheorem{proposition}{Proposition}
\newtheorem{lemma}{Lemma}
\newtheorem{corollary}{Corollary}
\newtheorem{principle}{Principle}
\newtheorem{grammarrule}{Rule}

\newcommand{\II}{\mathcal{I}}
\newcommand{\IIlip}{\mathcal{I}_{\mathrm{Lip}}}
\newcommand{\IImax}{\mathcal{I}_{\max}}
\newcommand{\Vsix}{V_{600}}
\newcommand{\Vtf}{V_{24}}
\newcommand{\Cphi}{C_\varphi}
\newcommand{\NH}{N_H}
\newcommand{\Q}{\mathbb{Q}}
\newcommand{\Z}{\mathbb{Z}}
\newcommand{\R}{\mathbb{R}}
\newcommand{\C}{\mathbb{C}}
\newcommand{\N}{\mathbb{N}}
\newcommand{\OK}{\mathcal{O}_K}
\newcommand{\Lambt}{\Lambda_{12}}
\newcommand{\Esix}{E_{600}}

\newtcolorbox{dictionarybox}[1][]{
  colback=blue!2,
  colframe=blue!40!black,
  fonttitle=\bfseries,
  title=Dictionary entry,
  #1
}

\newtcolorbox{grammarbox}[1][]{
  colback=green!2,
  colframe=green!40!black,
  fonttitle=\bfseries,
  title=Grammar rule,
  #1
}

\newtcolorbox{transbox}[1][]{
  colback=orange!2,
  colframe=orange!40!black,
  fonttitle=\bfseries,
  title=Translation,
  #1
}

\newtcolorbox{scopebox}[1][]{
  colback=gray!4,
  colframe=gray!50!black,
  fonttitle=\bfseries,
  title=Scope,
  #1
}

\title{\textbf{The Translation Layer}\\[0.4em]
\large Polytope Overlap as the Empirical Grammar of the $V_{600}$ Substrate, with Dictionary and Translation Rules}

\author{Lee Smart\\
\small Institute of Vibrational Field Dynamics\\
\small \texttt{contact@vibrationalfielddynamics.org}}

\date{2026-05-29 \quad (pre-peer-review preprint, v1.0.0-rc1)}

\begin{document}
\maketitle

\begin{abstract}
\noindent
The $V_{600}$ closure programme has been building a mathematical
substrate --- the icosian triad $(\II, \mathcal{G}, \Cphi)$ on the
$120$-vertex set of the regular $600$-cell --- and applying it
across mathematics, cosmology, biology and number theory. A natural
worry is that the substrate is functioning as a metaphor system
rather than as a working formal language. This paper makes the
case, with reproducible engine output, that the substrate now
behaves as a language in the technical sense: it has a
\emph{dictionary} (a typed map from substrate objects to candidate
interpretive referents), it has a \emph{grammar} (a finite set of
admissible transformations with explicit well-formedness
classes), and it has \emph{translation rules} (admissibility +
dictionary $\to$ candidate observable).

The grammar is not declarative. It is exhibited empirically by the
\emph{closure transformation engine}: a computational artefact that
computes, on $\Vsix$ and on polytope-overlap structures derived from
it, which transformations preserve closure and which break it.
Three load-bearing results pin the grammar down concretely:

\begin{enumerate}[leftmargin=*]
  \item The Schl\"afli decomposition
        $\Vsix = C_0 \sqcup C_1 \sqcup C_2 \sqcup C_3 \sqcup C_4$
        into five right cosets of the binary tetrahedral subgroup
        $\Vtf = 2T \subset 2I = \Vsix$, each isomorphic to the
        24-cell, with induced faithful $A_5$ action on the coset
        index --- five reproducible certificates.
  \item The $\lambda = 12$ \emph{lift channel}: every per-coset
        24-cell $\lambda = 12$ eigenvector extends to a full-$\Vsix$
        Laplacian $\lambda = 12$ eigenvector. The lifted space is
        $10$-dim inside $\dim \Esix(12) = 25$, with the residual
        $15$-dim block carrying a non-trivial $A_5$ character.
  \item The $\lambda = 12$ \emph{uniqueness scan}: across all tested
        subgroups $H \subset 2I$ and all global eigenvalues
        $\lambda'$, $\lambda' = 12$ is the only value at which the
        full lift channel opens. Every other $(H, \lambda')$ either
        produces a strictly smaller channel or no channel at all.
\end{enumerate}

These three results are the grammar: a positive rule
($\lambda = 12$ is admissible), a uniqueness claim ($\lambda \ne 12$
is forbidden), and a syntactic structure (the $A_5$ character split
of the lifted vs.\ residual blocks). On top of this grammar sit
two further classes of artefact: a \emph{transformation taxonomy}
that names twelve legal moves and five admissibility classes, with
which physical-interpretation labels are formally tied; and a
relational reasoning system (\texttt{closure\_ai} / \texttt{min\_aria})
that uses the engine's admissibility verdicts as its own state-
transition rules, demonstrating that the language is
\emph{compositional and productive} in the formal-language sense.

The paper's claim is not that this language settles physics. It is
that the substrate has crossed the threshold from metaphor to
formal language: dictionary, grammar, and translation are all
exhibited on reproducible artefacts. The substrate can be assessed
by the standards used to assess other working mathematical
languages.
\end{abstract}

\bigskip

\noindent\textbf{Status.} Pre-peer-review open research preprint.
Not independently validated. No claim of physical derivation, no
claim of consciousness reduction, no claim of RH / GRH settlement.
Every interpretive translation is explicitly a candidate
identification, conditional on the bridge hypotheses listed in
\S\ref{sec:scope}.

\bigskip

\tableofcontents

\newpage

% =====================================================================
\section{Introduction: the translation problem}
\label{sec:intro}

\subsection{Why ``language'' is the right shape, not just the right word}

The $V_{600}$ closure programme has been accumulating outputs of a
particular shape. Spectra of finite operators get identified with
families of physical observables. Cosets of finite groups get
identified with sectors of the Standard Model. Hilbert modular
$L$-functions get identified with classical zeta factors. The
recurrence is striking and the underlying object --- the icosian
triad $(\II, \mathcal{G}, \Cphi)$ on $\Vsix$ --- keeps doing the same
kind of work in each case. The natural diagnosis is that the
substrate is functioning as a language: a fixed vocabulary plus a
set of composition rules from which a wide range of declarative
sentences can be formed.

This paper takes that diagnosis literally and asks what would have
to be true for it to be correct. A working formal language requires
three components:

\begin{enumerate}[leftmargin=*]
  \item A \emph{dictionary} (lexicon): a typed map from primitive
        substrate objects to candidate referents in the
        interpretive layer.
  \item A \emph{grammar}: a finite set of rules saying which
        compositions of primitives are well-formed, which are
        forbidden, and which are conditionally admissible.
  \item A \emph{translation procedure}: a rule that converts a
        well-formed sentence of the substrate language into a
        statement in the interpretive layer.
\end{enumerate}

Without (i) the substrate is a graph. Without (ii) the substrate is
a metaphor system. Without (iii) the substrate is a
dictionary-plus-grammar with no semantics. \emph{All three} are
required.

\subsection{What is concrete and publishable here}

We claim, and exhibit, all three on reproducible engine output. The
key empirical artefact is the \emph{closure transformation engine}.
It is a computational pipeline that, given a graph or polytope
specification, computes:

\begin{itemize}[leftmargin=*]
  \item closure validation (does the structure close consistently?);
  \item invariant extraction (Euler characteristic, spectra,
        curvature class);
  \item duality and projection operators;
  \item admissibility verdicts for proposed transformations.
\end{itemize}

The engine does not assume physical content. It is a
\emph{mathematical} pipeline. Its physical-interpretation labels
are deliberately conservative; every label is keyed off an abstract
\texttt{Admissibility\-Class} and is phrased as a candidate
identification.

The reason the engine carries the paper is that it lets us write
down the grammar of the substrate \emph{empirically}, by observing
which transformations preserve closure and which break it. The
$\lambda = 12$ lift channel and its uniqueness across alternative
subgroups and eigenvalues is the cleanest example: an exhibited
\emph{positive} grammar rule plus an exhibited \emph{negative}
constraint, both reproducible.

\subsection{What this paper is and is not}

\begin{scopebox}
\textbf{Is.} An empirical demonstration that the $V_{600}$
substrate now meets the three formal criteria for a language ---
dictionary, grammar, translation --- with reproducible engine
output as the load-bearing evidence.

\textbf{Is not.} A claim that the language settles physics, life,
mind, or RH. Every interpretive translation in this paper is a
candidate identification, conditional on a named bridge
hypothesis.
\end{scopebox}

\subsection{Roadmap}

\S\ref{sec:layers} states the three-layer framing.
\S\ref{sec:engine} describes the closure transformation engine and
its outputs.
\S\ref{sec:dictionary} exhibits the dictionary.
\S\ref{sec:grammar} exhibits the grammar from polytope overlap,
with the Schl\"afli decomposition, the $\lambda = 12$ lift, and the
$\lambda = 12$ uniqueness scan as the three load-bearing
sub-results.
\S\ref{sec:translation} exhibits the translation procedure.
\S\ref{sec:scope} states the substrate's expressive scope ---
explicitly, what it cannot say.
\S\ref{sec:compositional} shows compositionality by lifting the
grammar to the relational-AI architecture \texttt{closure\_ai} /
\texttt{min\_aria}.
\S\ref{sec:comparison} compares the substrate to other working
mathematical languages (group theory, twistor theory,
amplituhedron, LQG).
\S\ref{sec:openness} states open translation problems.
\S\ref{sec:conclusion} concludes.

% =====================================================================
\section{Three layers: dictionary, grammar, interpretation}
\label{sec:layers}

The three-layer framing originated in our evidence-ledger audit
\cite{evidenceLedger} as a way to keep ``$\Vsix$ appears
everywhere'' from quietly counting as evidence. Here we lift the
same framing into a positive constructive frame.

\subsection{Layer 1: dictionary}

A dictionary entry is a typed map
$d_i : \Sigma_i \to \mathrm{Ref}_i$ that takes a substrate object
$\Sigma_i$ (a vertex, a coset, an eigenspace, a character) to a
candidate referent. A dictionary entry is \emph{not evidence}: it
is a translation \emph{possibility}, awaiting confirmation by
grammar.

\subsection{Layer 2: grammar}

A grammar rule is a deterministic statement about which
compositions of substrate objects are admissible. The closure
transformation engine produces grammar rules
\emph{operationally}: a transformation $T$ applied to a closed
state $S$ either preserves closure (\texttt{Admissibility = EXACT}
or \texttt{CANDIDATE}) or breaks it
(\texttt{APPROXIMATE} / \texttt{BROKEN} / \texttt{DEGENERATE}).
The set of $(T, S)$ pairs returning \texttt{EXACT} is the grammar.

\subsection{Layer 3: translation}

Translation maps a dictionary entry $d_i$ plus a grammar verdict
$v$ into an interpretive claim
$\mathrm{Ref}_i \mid v$. Crucially: \emph{translation requires
grammar.} A dictionary entry $\Vsix \leadsto \text{``substrate of
quantum coherence''}$ is not evidence until the grammar verdict on
the relevant composition has been computed. The engine computes
the verdict; the translation is whatever the verdict permits.

\begin{tcolorbox}[colback=gray!5,colframe=gray!50!black,
  title={The three layers, restated}]
\begin{tabular}{lll}
\toprule
\textbf{Layer} & \textbf{What it is} & \textbf{Evidential status} \\
\midrule
1 Dictionary  & typed substrate $\to$ referent map
              & possibility (not evidence) \\
2 Grammar     & engine-computed admissibility verdict
              & empirical (theorem-grade or reproducible) \\
3 Translation & dictionary + grammar $\to$ interpretive claim
              & conditional on grammar + bridge hyp.\ \\
\bottomrule
\end{tabular}
\end{tcolorbox}

A common failure mode in informal cross-disciplinary work is to
present a Layer-1 dictionary as if it were a Layer-3 translation.
The audit pattern of this paper is the opposite: every Layer-3
claim points to a Layer-2 verdict and a Layer-1 entry, and the
reader can disagree at any of the three.

% =====================================================================
\section{The closure transformation engine}
\label{sec:engine}

\subsection{What the engine is}

The closure transformation engine \cite{transformationEngine} is a
computational pipeline. Its inputs are graph specifications and
polytope overlap structures derived from $\Vsix$. Its outputs are
closure verdicts, invariant reports, transformation taxonomy
records, admissibility classes and (when requested) candidate
physical-interpretation labels.

The engine has three core modules:

\begin{description}[leftmargin=*]
  \item[\texttt{core/}] closure validation
        (\texttt{closure\_validator.py}), invariant extraction
        (\texttt{invariants.py}, \texttt{spectrum.py},
        \texttt{curvature.py}), duality
        (\texttt{duality.py}), projection
        (\texttt{projection.py}), residuals
        (\texttt{residuals.py}), and Schl\"afli-symbol intake
        (\texttt{schlafli.py}).
  \item[\texttt{transformations/}] twelve transformation kinds
        (DUAL, SHELL\_PROJECTION, BOUNDARY\_PROJECTION,
        SPECTRAL\_PROJECTION, CYCLE\_PROJECTION, QUOTIENT,
        SUBGROUP\_RESTRICTION, EDGE\_CONTRACTION,
        VERTEX\_DELETION, EDGE\_DELETION, DIMENSIONAL\_LIFT,
        DIMENSIONAL\_PROJECTION). Each carries a typed
        \texttt{TransformationOutcome} record.
  \item[\texttt{physics/}] admissibility classification
        (\texttt{admissibility.py}), transition rules between
        decorated states (\texttt{transition\_rules.py}), the
        24-600 spectral bridge (\texttt{bridge\_24\_600.py}), the
        $\lambda = 12$ channel
        (\texttt{lambda12\_a5\_channel.py}), the universality scan
        (\texttt{universality\_lambda12.py}), shell hierarchy and
        gauge classification, and the (deliberately conservative)
        \texttt{physical\_interpretation.py} label table.
\end{description}

\subsection{Why the engine is the load-bearing artefact}

The engine matters here because it is the place where the
substrate's grammar is no longer a verbal claim. Every grammar
rule we state in \S\ref{sec:grammar} is reproducible from one of
the engine's outputs. Three artefacts in particular do most of the
work:

\begin{itemize}[leftmargin=*]
  \item the Schl\"afli decomposition module
        (\texttt{decomposition.py}), with five exact certificates;
  \item the $\lambda = 12$ spectral-bridge module
        (\texttt{keystone.py}), with the lift theorem and the
        $A_5$ character decomposition;
  \item the $\lambda = 12$ universality module
        (\texttt{universality\_lambda12.py}), which scans all
        subgroups $H \subset 2I$ and all $\Vsix$ Laplacian
        eigenvalues to test whether $\lambda = 12$ is the unique
        full-lift channel.
\end{itemize}

These three artefacts together yield the
positive-and-negative grammar rule discussed in
\S\ref{sec:grammar:lambda12}.

% =====================================================================
\section{The dictionary}
\label{sec:dictionary}

This section exhibits the dictionary. Each entry maps a substrate
object to one or more candidate interpretive referents. The
referents are typed --- ``mathematical'', ``physical (candidate)'',
``cognitive (candidate)'', ``ontological (interpretive)'' --- and
each entry carries an evidential weight inherited from its
grammar-layer status.

\subsection{Primitives}

\begin{dictionarybox}
\textbf{$\Vsix = 2I$.} 120-element binary icosahedral group,
realised as 120 unit icosian quaternions.
\\[0.3em]
\textbf{Mathematical referent:} vertex set of regular $600$-cell;
finite reflection-group orbit under $W(H_4)$.\\
\textbf{Candidate physical referent:} closure-substrate carrier.\\
\textbf{Evidential weight:} Layer-1; weight comes from grammar.
\end{dictionarybox}

\begin{dictionarybox}
\textbf{$\Vtf = 2T \subset \Vsix$.} 24-element binary tetrahedral
subgroup, realised as the $\sigma$-fixed sub-lattice of $\Vsix$.
\\[0.3em]
\textbf{Mathematical referent:} 24-cell vertex set; the
fundamental coset block in the Schl\"afli decomposition.\\
\textbf{Candidate physical referent:} closure ``unit cell'' /
fundamental sector.\\
\textbf{Evidential weight:} Layer-1.
\end{dictionarybox}

\begin{dictionarybox}
\textbf{$A_1$ (edge-adjacency operator on $\Vsix$).} $120 \times 120$
matrix with $(A_1)_{ij} = 1$ iff $v_i, v_j$ are vertex-adjacent in
$\Vsix$.
\\[0.3em]
\textbf{Mathematical referent:} graph-Laplacian generator; defines
$L = 12 \cdot I - A_1$.\\
\textbf{Spectrum:} 9 eigenvalues in $\Z[\varphi]$ with
multiplicities $(1, 4, 9, 16, 25, 36, 9, 16, 4)$ summing to $120$.\\
\textbf{Evidential weight:} Layer-2 (theorem-grade).
\end{dictionarybox}

\begin{dictionarybox}
\textbf{$\Cphi = L + \varphi^{-2} I = (12 + \varphi^{-2}) I - A_1$.}
The closure operator.
\\[0.3em]
\textbf{Mathematical referent:} self-adjoint positive operator
inducing the $94 + 13 + 13$ $\sigma$-pairing decomposition.\\
\textbf{Candidate referent:} the ``closure verb'' --- the
mathematical instantiation of ``$X$ exists $\iff
\mathcal{C}(X) = X$''.\\
\textbf{Evidential weight:} Layer-2 (theorem-grade) for the math;
Layer-3 (interpretation) for the existence reading.
\end{dictionarybox}

\begin{dictionarybox}
\textbf{$\NH : \II \to \Z[\varphi]$,
$\NH(q) = q_0^2 + q_1^2 + q_2^2 + q_3^2$.} The generation
operator.
\\[0.3em]
\textbf{Mathematical referent:} quaternary norm form, universal
for totally positive $\Z[\varphi]$-primes by Eichler--Hijikata.\\
\textbf{Evidential weight:} Layer-2 (theorem-grade).
\end{dictionarybox}

\subsection{Compound primitives from polytope overlap}

\begin{dictionarybox}
\textbf{The five 24-cell cosets $C_0, \ldots, C_4$.}
$\Vsix = \bigsqcup_{i=0}^4 C_i$ with $|C_i| = 24$, each $C_i$
isomorphic to the 24-cell.
\\[0.3em]
\textbf{Mathematical referent:} right cosets of $\Vtf$ in $\Vsix$;
five copies of the 24-cell tiled into one 600-cell.\\
\textbf{Candidate physical referent:} the five \emph{building
blocks} of any physical sector built on $\Vsix$; the carriers of
the $A_5$ representation.\\
\textbf{Evidential weight:} Layer-2 (theorem-grade) for the
decomposition; Layer-3 for ``building blocks''.
\end{dictionarybox}

\begin{dictionarybox}
\textbf{$\Esix(\lambda)$.} The $\lambda$-eigenspace of the
$\Vsix$ adjacency Laplacian.
\\[0.3em]
\textbf{Mathematical referent:} a finite-dimensional $W(H_4)$-rep.\\
\textbf{Distinguished entry:} $\dim \Esix(12) = 25$.\\
\textbf{Candidate physical referent:} candidate
spectral-mode-family. The labelling per $\lambda$ is in
\cite{icosianTriad}.
\end{dictionarybox}

\begin{dictionarybox}
\textbf{The $\lambda = 12$ lift channel $\Lambt$.} The
$10$-dimensional subspace of $\Esix(12)$ obtained by zero-extending
the per-coset $\lambda = 12$ 24-cell eigenvectors.
\\[0.3em]
\textbf{Mathematical referent:} \emph{the carrier of the grammar
rule of \S\ref{sec:grammar:lambda12}.} Its complement
$\Esix(12) \setminus \Lambt$ is the 15-dim residual block, with a
non-trivial $A_5$ character.\\
\textbf{Evidential weight:} Layer-2 (theorem-grade).
\end{dictionarybox}

\subsection{Dictionary table}

The full machine-readable dictionary is at
\texttt{data/dictionary.csv}. It records, per row, the substrate
object, the mathematical referent, the candidate physical
referent, the grammar-layer status, the bridge hypothesis (if
any), and the evidence-ledger row of \cite{evidenceLedger} that
audits the entry.

% =====================================================================
\section{Polytope overlap as grammar}
\label{sec:grammar}

This is the load-bearing section of the paper. We exhibit three
grammar rules derived empirically from the closure transformation
engine. Two are positive (the engine returns
\texttt{EXACT}); one is negative (the engine returns no channel).

\subsection{The Schl\"afli decomposition as the grammar's primitive
sentence}
\label{sec:grammar:schlafli}

\begin{theorem}[Schl\"afli decomposition; engine certificate]
\label{thm:schlafli}
The 120-vertex set $\Vsix = 2I$ admits the exact partition
\[
  \Vsix \;=\; C_0 \;\sqcup\; C_1 \;\sqcup\; C_2 \;\sqcup\; C_3 \;\sqcup\; C_4
\]
where $C_0 = \Vtf = 2T$ (the $\sigma$-fixed sub-lattice) and
$C_1, \ldots, C_4$ are the four $\Vtf$-left-multiplication orbits
on the 96 $\sigma$-mobile vertices. Each $C_i$ is exactly a coset of
$\Vtf$, has cardinality 24, and is isometric to the 24-cell vertex
set with d=1 graph $8$-regular on 96 edges. The induced action of
$\Vsix$ on the coset set $\{C_0, \ldots, C_4\}$ descends through
$\{\pm 1\}$ to a faithful transitive action of $A_5$ on five
points.
\end{theorem}

The theorem is verified by five reproducible certificates from
\texttt{closure\_transform\_engine.decomposition}: \texttt{V\_600}
construction, \texttt{V\_24} subgroup, right cosets, each coset
is a 24-cell, $A_5$ coset action. All five are exact over
$\Q(\sqrt 5)$.

\paragraph{Why this is a grammar primitive.} Because every
subsequent grammar rule in this paper either references the five
cosets or references the $A_5$ representation they carry.
Theorem~\ref{thm:schlafli} is the minimal compositional unit on
which the rest of the grammar is built.

\subsection{The $\lambda = 12$ lift rule}
\label{sec:grammar:lambda12}

\begin{theorem}[Per-coset $\lambda = 12$ lift theorem; engine
certificate]
\label{thm:lambda12}
Let $L_{600}$ be the d=1 graph Laplacian on $\Vsix$
($12$-regular, $720$ edges). For each $i \in \{0,1,2,3,4\}$ let
$L_{C_i}$ be the d=1 graph Laplacian on the coset $C_i$
($8$-regular, $96$ edges, isomorphic to the 24-cell edge graph).
For each $i$, the $\lambda = 12$ eigenspace $E_{C_i}(12)$ of
$L_{C_i}$ has dimension 2. The zero-extension of any
$v \in E_{C_i}(12)$ to $\R^{120}$ is an eigenvector of $L_{600}$
at eigenvalue exactly $12$. Summing over $i$ yields a
$10$-dimensional subspace $\Lambt \subset \Esix(12)$. The
complementary 15-dimensional residual block carries a non-trivial
$A_5$ character orthogonal to the lifted block.
\end{theorem}

The theorem is verified by
\texttt{closure\_transform\_engine.keystone.exact\_certify\_lambda12\_lift}
and the $A_5$ character calculation
\texttt{exact\_a5\_characters}.

\begin{grammarbox}
\textbf{Rule (positive).} ``Lift along the Schl\"afli decomposition
at $\lambda = 12$ is admissible.''
\\
\emph{Substrate side:} per-coset 24-cell eigenvector at
$\lambda = 12$.
\\
\emph{Grammar verdict:} the engine returns
\texttt{Admissibility = EXACT} with $\|L_{600} v - 12 v\| <
\epsilon_{\mathrm{mach}}$ for every input.
\\
\emph{Composition:} the lift is linear and commutes with the
$A_5$ action on cosets.
\end{grammarbox}

\subsection{The $\lambda = 12$ uniqueness scan: a forbidden-move
rule}
\label{sec:grammar:uniqueness}

The positive rule of \S\ref{sec:grammar:lambda12} is not yet a
grammar; a grammar must also tell us which moves are
\emph{forbidden}. The empirical answer is supplied by the
universality scan in
\texttt{closure\_transform\_engine.physics.universality\_lambda12}.

\paragraph{Setup.} For each subgroup $H \subset 2I$ in a finite
test set and each eigenvalue $\lambda' \in \mathrm{spec}(L_{600})$,
the scan computes
\[
\mathrm{Lift}(H, \lambda') \;=\; \bigoplus_i \big\{ v \in \R^{V_{600}} :
 \mathrm{supp}(v) \subseteq C_i^{(H)},\;\; L_{600} v = \lambda' v \big\},
\]
where $C_i^{(H)}$ ranges over the right cosets of $H$ in $2I$. The
scan reports the dimension of $\mathrm{Lift}(H, \lambda')$ for every
pair $(H, \lambda')$, and flags every pair whose channel dimension
is non-zero.

\paragraph{Empirical finding.} Across the tested subgroup set, the
\emph{full} lift channel
($\mathrm{Lift}(H, \lambda') = \dim \Esix(\lambda')$ from per-coset
support alone) opens exactly when $H = \Vtf$ and $\lambda' = 12$.
Every other pair returns either a strictly smaller channel or no
channel at all. The five negative controls
$\lambda' \in \{0, 4, 8, 10\}$ all return zero per-coset lift into
$L_{600}$.

\begin{grammarbox}
\textbf{Rule (negative).} ``Polytope-overlap lift at
$\lambda' \ne 12$ is forbidden.''
\\
\emph{Substrate side:} per-$H$-coset eigenvector at $\lambda'$.
\\
\emph{Grammar verdict:} for $\lambda' \in \{0, 4, 8, 10\}$ the
engine returns \texttt{Admissibility = BROKEN}.
\\
\emph{Composition:} the forbidden-move rule survives change of
subgroup $H$ in the tested set.
\end{grammarbox}

The conjunction ``$\lambda = 12$ admissible, all others
forbidden'' is the first grammar rule the substrate has at the
sentence level. It is a clean empirical \emph{constraint}, not a
labelling choice.

\subsection{Transformation kinds as the legal-move set}
\label{sec:grammar:moves}

The engine's transformation taxonomy provides the substrate's
operational verbs: the twelve transformation kinds (DUAL,
SHELL\_PROJECTION, BOUNDARY\_PROJECTION, SPECTRAL\_PROJECTION,
CYCLE\_PROJECTION, QUOTIENT, SUBGROUP\_RESTRICTION,
EDGE\_CONTRACTION, VERTEX\_DELETION, EDGE\_DELETION,
DIMENSIONAL\_LIFT, DIMENSIONAL\_PROJECTION). Each carries an
explicit input typing, an output typing, and a closure-residual
recomputation. A move is \emph{legal} when its admissibility class
returns \texttt{EXACT} or \texttt{CANDIDATE}; it is \emph{illegal}
when it returns \texttt{BROKEN} or \texttt{DEGENERATE}.

This is the substrate's grammar in a strict computational sense.
The full move table is at \texttt{data/grammar.csv}, together with
the per-move pre-and-post invariant lists.

\subsection{Admissibility classes as well-formedness}
\label{sec:grammar:wellformed}

The engine resolves every transformation into one of five
admissibility classes:

\begin{itemize}[leftmargin=*]
  \item \texttt{EXACT} --- closure preserved within numerical
        tolerance.
  \item \texttt{CANDIDATE} --- near-closed, low residual.
  \item \texttt{APPROXIMATE} --- noticeable residual but
        interpretable.
  \item \texttt{BROKEN} --- closure lost.
  \item \texttt{DEGENERATE} --- transformation produced a trivial
        or empty target.
\end{itemize}

A well-formed sentence of the substrate language is, by
definition, a chain of moves each of which terminates in
\texttt{EXACT} or \texttt{CANDIDATE}. The substrate's grammar is
this well-formedness predicate.

\subsection{Transition rules as composition}
\label{sec:grammar:transitions}

The engine's \texttt{TransitionRule} record specifies the
default-allowed transitions between decorated closure states:
\texttt{preserve\_topology}, \texttt{preserve\_charge\_like\_labels},
\texttt{preserve\_boundary\_closure}, with
\texttt{max\_residual\_delta = 0.01} and optional
\texttt{allow\_phase\_defect}. Each transition is resolved to a
verdict from the set \{allowed, forbidden, allowed-as-virtual,
allowed-as-decay-candidate, requires-boundary-change\}. This is
the substrate's compositional rule: how shorter well-formed
sentences combine into longer ones.

% =====================================================================
\section{Translation rules}
\label{sec:translation}

Translation is the operation $T : (d_i, v) \mapsto r_i \mid v$.
The dictionary supplies $d_i$; the grammar supplies the verdict
$v$; translation produces a candidate interpretive claim
$r_i \mid v$.

\subsection{The translation procedure, formally}

\begin{principle}[Translation]
\label{prin:translation}
A translation from substrate object $\Sigma_i$ to candidate
referent $\mathrm{Ref}_i$ is admissible iff:

\begin{enumerate}[leftmargin=*,label=(\roman*)]
  \item A dictionary entry $d_i : \Sigma_i \to \mathrm{Ref}_i$
        exists in \texttt{data/dictionary.csv}.
  \item There is a sentence $s$ of the substrate language using
        $\Sigma_i$ as a primitive whose grammar verdict $v(s)$ is
        \texttt{EXACT} or \texttt{CANDIDATE}.
  \item A named bridge hypothesis $H_{i,\mathrm{Ref}}$ is supplied,
        connecting $\Sigma_i \mid v(s)$ to the interpretive
        domain.
\end{enumerate}

The translation is then $r_i \mid (v(s) \wedge H_{i,\mathrm{Ref}})$.
\end{principle}

The principle ties Layer 3 to Layer 2 to Layer 1. Without (ii) the
translation is metaphor. Without (iii) the translation is a Hilbert
space without an inner product: well-formed, semantically empty.

\subsection{Worked translation: closed admissibility $\to$ stable
observable candidate}

\begin{transbox}
\textbf{Substrate sentence.} A decorated closure state $S$ over
$\Vsix$, with the $\lambda = 12$ lift channel populated and zero
boundary defect.
\\
\textbf{Grammar verdict.} \texttt{Admissibility = CLOSED} (per
\texttt{physical\_interpretation.py}, line \texttt{\_INTERPRETATIONS\!
[AdmissibilityClass.CLOSED]}).
\\
\textbf{Dictionary entries used.} $\Vsix$ as
closure-substrate; $\Lambt$ as carrier of the closure mode-family.
\\
\textbf{Bridge hypothesis required.} H-Observable: stable
mathematical closure corresponds to a candidate physical
observable.
\\
\textbf{Translation.} ``The closure-substrate carries a candidate
stable observable, conditional on H-Observable.''
\\
\textbf{Audit caveat (engine-enforced).} ``Not yet a physical
derivation; closure does not imply realisability.''
\end{transbox}

This is the deliberate conservatism of the engine reflected at the
paper level. The engine attaches an explicit caution to every
physical-interpretation label and the translation principle
inherits it.

\subsection{Worked translation: near-closed $\to$ resonance
candidate}

\begin{transbox}
\textbf{Substrate sentence.} Decorated state $S'$ with small but
non-zero closure residual on the boundary.
\\
\textbf{Grammar verdict.} \texttt{Admissibility = NEAR\_CLOSED}.
\\
\textbf{Dictionary entries.} same as above plus a boundary-decoration
field.
\\
\textbf{Bridge hypothesis.} H-Resonance: near-closure
$\leftrightarrow$ candidate resonance / transient bound state.
\\
\textbf{Translation.} ``Candidate resonance or transient bound
state, conditional on H-Resonance and on the boundary decoration
not being an artefact.''
\\
\textbf{Audit caveat.} ``May be an artefact of the chosen
decoration --- try alternative phase or boundary assignments
before drawing conclusions.''
\end{transbox}

\subsection{Worked translation: $\lambda = 12$ universality $\to$
substrate mode-family uniqueness}

\begin{transbox}
\textbf{Substrate sentence.} The $\lambda = 12$ lift theorem and
its uniqueness scan.
\\
\textbf{Grammar verdict.} \texttt{EXACT} for the lift theorem;
\texttt{BROKEN} or smaller-channel for every other
$(H, \lambda')$ in the tested set.
\\
\textbf{Dictionary entries.} $\Vsix$, the five cosets, $\Lambt$.
\\
\textbf{Bridge hypothesis.} H-ModeUnique: substrate mode-family
selection corresponds to grammar-rule selection.
\\
\textbf{Translation.} ``The substrate selects exactly one
mode-family across the tested choice of subgroup and eigenvalue,
conditional on H-ModeUnique.''
\\
\textbf{Audit caveat.} The test set of subgroups is finite. Extending
the test set is the open work.
\end{transbox}

\subsection{What translation is \emph{not}}

\begin{itemize}[leftmargin=*]
  \item Translation is not a physical derivation. A translation
        produces a \emph{candidate} identification. The candidate
        becomes a derivation only when the bridge hypothesis is
        independently closed.
  \item Translation is not equivalence. The interpretive referent
        may have structure the substrate cannot represent.
        Translation is one-way unless the dictionary entry has an
        inverse.
  \item Translation is not falsification. A failing translation
        does not falsify the substrate; it falsifies the
        translation or the bridge hypothesis. The substrate
        remains.
\end{itemize}

% =====================================================================
\section{Expressive scope: what the substrate can and cannot say}
\label{sec:scope}

A working formal language has an \emph{expressive scope}. There are
sentences inside it and sentences it cannot form. The grammar rules
of \S\ref{sec:grammar} let us state the substrate's scope
explicitly.

\subsection{What the substrate can say}

\begin{enumerate}[leftmargin=*]
  \item Spectra of self-adjoint operators on $\Vsix$, with
        eigenvalues in $\Z[\varphi]$.
  \item Subspaces invariant under $W(H_4)$, decomposed into
        irreducibles of the Coxeter group.
  \item Cosets of finite subgroups $H \subset 2I$ and the lift
        channels they open.
  \item $\sigma$-pairing decompositions of subspaces (the
        $94 + 13 + 13$ split).
  \item Quaternary norm forms $\NH$ and the totally positive
        $\Z[\varphi]$-primes they enumerate.
  \item Hilbert modular $L$-functions over $K = \Q(\sqrt 5)$
        factoring through the substrate theta series.
  \item Composition of any of the above under transformation
        kinds in the move table.
\end{enumerate}

\subsection{What the substrate cannot say}

\begin{enumerate}[leftmargin=*]
  \item Continuous deformations of $\Vsix$. The substrate is
        finite. Smooth deformations have to be carried by the
        interpretive layer.
  \item Statements that are not $W(H_4)$-covariant. The grammar
        forbids them. (This is the substrate's first explicit
        \emph{negation}.)
  \item Statements that require a subgroup outside the test set
        of the universality scan. These remain open.
  \item Statements about objects outside $\Q(\sqrt 5)$. Other
        biquadratic fields are reachable only via the
        \texttt{DIMENSIONAL\_LIFT} kind, and only with named
        bridge hypotheses.
  \item Anything in the interpretive layer that has no dictionary
        entry. Without a dictionary entry, there is no
        translation.
\end{enumerate}

\subsection{Bridge hypotheses currently consumed by translation}

\begin{description}[leftmargin=*]
  \item[H-Observable.] Closure $\leftrightarrow$ candidate
        observable. Open.
  \item[H-Resonance.] Near-closure $\leftrightarrow$ candidate
        resonance. Open.
  \item[H-ModeUnique.] $\lambda = 12$ universality
        $\leftrightarrow$ substrate mode-family uniqueness. Open.
  \item[H-CortexClosure.] Cortical phase closure
        $\leftrightarrow$ closure observable
        \cite{closurePicture}. Open.
  \item[H-MT-Bridge.] Microtubule lattice $\leftrightarrow$
        candidate substrate carrier. Open (exploratory).
\end{description}

The four-tier scope discipline of
\cite[\S 9]{closurePicture} and the constraint-forcing criterion
of \cite[\S 5]{evidenceLedger} both apply unchanged here.

% =====================================================================
\section{Compositionality: the language is productive}
\label{sec:compositional}

A working language is more than dictionary + grammar + translation:
it is \emph{compositional} and \emph{productive}. Finitely many
primitives plus a grammar yield infinitely many well-formed
sentences. We show that the substrate has this property by
exhibiting a reasoning system built on top of the engine.

\subsection{\texttt{closure\_ai}: relational reasoning over
substrate states}

\texttt{closure\_ai} \cite{closureAI} is a relational reasoning
architecture whose hidden state is a decorated cognitive graph,
whose reasoning steps are graph transformations from the engine,
and whose attention is invariant-aware projection. Crucially:
\emph{outputs are emitted only after the admissibility gate
accepts the final state.}

\paragraph{Why this matters here.} It shows that the engine's
admissibility classes are usable as a complete state-transition
predicate. The system runs because the grammar of
\S\ref{sec:grammar:wellformed} is decidable. A grammar that
classified well-formedness ambiguously could not drive a
reasoning loop. The fact that \texttt{closure\_ai} runs is
empirical evidence that the substrate grammar is a
\emph{functioning} formal language, not just a labelled set of
moves.

\subsection{\texttt{min\_aria}: vector-state substrate carrier}

\texttt{min\_aria} \cite{minAria} is a minimal substrate-aware
agent carrying a 7-dimensional vector state
\[
S_t \;=\; [
  \mathrm{identity\_stability},\
  \mathrm{memory\_load},\
  \mathrm{novelty\_score},\
  \mathrm{closure\_residual},
\]
\[\hphantom{S_t \;=\; [}
  \mathrm{perturbation\_level},\
  \mathrm{governance\_integrity},\
  \mathrm{replay\_integrity}].
\]
The closure-residual component is the engine's well-formedness
predicate of \S\ref{sec:grammar:wellformed}; the perturbation and
governance components are direct readouts of admissibility-class
transitions. The agent maintains identity under hostile prompts
($\mathrm{identity\_drift} \approx 0$) and recovers closure under
perturbation ($\mathrm{closure\_residual} \to 0$). The numerical
proofs of these claims live in
\texttt{min\_aria/test\_min\_aria\_v2.py}.

\paragraph{What this proves.} The substrate's grammar is
compositional enough to drive a stateful agent. That is not a
weak claim. It is the property that distinguishes ``a labelled
algebraic structure'' from ``a formal language''.

\subsection{Productivity: finite vocabulary, unbounded
expressions}

The dictionary has finite size. The grammar has finite move set.
But the set of well-formed sentences is unbounded: composition of
transformations can produce decorated states of arbitrarily large
graph depth, and the engine's residual budget gives the well-formed
ones a clean stopping criterion. This is the precise
property linguistic theory calls \emph{productivity}.

\subsection{Falsifiability: the grammar refuses sentences}

The grammar refuses sentences. \texttt{BROKEN} verdicts are
returned every time a transformation breaks closure beyond
$\Delta_{\mathrm{max}} = 0.01$ in the residual budget. This is
exactly the property linguistic theory calls
\emph{discriminative power} --- the grammar produces ill-formed
sentences as well as well-formed ones.

\bigskip

\noindent
\textbf{Together:} dictionary $+$ decidable grammar $+$
productivity $+$ discriminative power are the four classical
criteria for a formal language. The substrate exhibits all four,
reproducibly, on engine artefacts.

% =====================================================================
\section{Comparison to other mathematical languages}
\label{sec:comparison}

The $V_{600}$ substrate is not the first mathematical language to
be exhibited in roughly this form. The comparison is informative
because it sets a recognisable scale.

\subsection{Group-representation theory}

The dictionary is the character table; the grammar is induced and
restricted representations and tensor products; the translation
layer is the assignment of irreducible representations to physical
multiplets. Productivity: a finite group generates an unbounded
sequence of well-formed compound representations. Falsifiability:
representation theory refuses tensor products that fail the
character orthogonality relations.

\subsection{Lie algebra and root systems}

The dictionary is roots and weights; the grammar is the structure
constants and Weyl group action; the translation layer is the
identification of simple roots with gauge generators. Productivity
is again unbounded; falsifiability is the Jacobi identity.

\subsection{Twistor theory}

The dictionary maps twistor space coordinates to null directions
in Minkowski space; the grammar is the Penrose transform; the
translation produces scattering-amplitude representations. The
twistor programme has explicit Layer-2 (grammar) results, Layer-3
(translation) results conditional on bridge hypotheses (the
Penrose transform's analyticity assumptions), and Layer-1
(dictionary) entries used across the programme.

\subsection{Amplituhedron programme}

The dictionary is positive Grassmannian cells; the grammar is the
boundary and triangulation calculus; the translation produces
$\mathcal{N}=4$ super-Yang--Mills amplitudes. The translation
\emph{is} the amplitude, so the bridge hypothesis is exceptionally
strong (rather than weak) in this case. The amplituhedron is the
nearest comparison to what the $V_{600}$ substrate aspires to be.

\subsection{Loop quantum gravity}

The dictionary is spin networks; the grammar is gauge-invariant
moves on the graph; the translation produces candidate
gravitational observables. The bridge hypotheses are explicit and
open.

\subsection{Where $V_{600}$ sits}

The $V_{600}$ substrate sits in this same shape but with three
distinguishing features:

\begin{itemize}[leftmargin=*]
  \item Its grammar is exhibited on a single \emph{fixed} finite
        substrate (the $600$-cell vertex set) rather than on a
        family of generators.
  \item It explicitly carries an evidence ledger
        \cite{evidenceLedger}, separating dictionary entries from
        constraint-forcing entries.
  \item It is, currently, the work of a single author working
        outside an institutional setting, released as pre-peer-
        review preprints. We are explicit that the comparison to
        the four programmes above is structural, not equal in
        stature.
\end{itemize}

% =====================================================================
\section{Open translation problems}
\label{sec:openness}

The translation layer is the youngest of the three layers and has
the most open work. We list the substantive open problems so that
future work has a focus.

\begin{enumerate}[leftmargin=*]
  \item \textbf{Subgroup test-set completion.} The
        $\lambda = 12$ universality scan
        (\texttt{universality\_lambda12.py}) was run over a finite
        test set of subgroups $H \subset 2I$. Extending the scan
        to \emph{all} subgroups of $2I$ would convert the negative
        rule of \S\ref{sec:grammar:uniqueness} from
        empirical-on-test-set into theorem-grade.
  \item \textbf{Bridge-hypothesis closure.}
        H-Observable, H-Resonance, H-ModeUnique,
        H-CortexClosure, H-MT-Bridge are all open. Each
        translation in \S\ref{sec:translation} stands or falls
        with its bridge.
  \item \textbf{Inverse-translation construction.} Every
        translation in this paper is one-way (substrate $\to$
        interpretation). Whether the dictionary has inverses ---
        whether every well-formed interpretive sentence has a
        substrate translation --- is open.
  \item \textbf{Compositionality at scale.} \texttt{closure\_ai}
        and \texttt{min\_aria} demonstrate compositionality at
        small scale. Whether the grammar remains decidable as the
        sentence length grows is an open empirical question; the
        residual budget gives the right place to look.
  \item \textbf{Bidirectional translation table.} The
        \texttt{data/translation.csv} CSV gives the current
        forward translation table. A reverse table
        (interpretation $\to$ candidate substrate sentence) does
        not yet exist; it is open work.
\end{enumerate}

% =====================================================================
\section{Conclusion}
\label{sec:conclusion}

This paper has argued that the $V_{600}$ substrate has crossed the
threshold from metaphor to formal language. The argument is not
declarative; it is empirical. The closure transformation engine
supplies the grammar, reproducibly. The Schl\"afli decomposition
and the $\lambda = 12$ lift-channel theorem are positive rules.
The $\lambda = 12$ uniqueness scan is a negative rule. The
transformation taxonomy supplies the move set. The admissibility
classes supply the well-formedness predicate. The translation
principle ties dictionary to grammar to interpretation. And
\texttt{closure\_ai} / \texttt{min\_aria} show that the grammar is
decidable enough to drive a reasoning system.

The substrate now meets the four classical criteria of a formal
language: typed dictionary, decidable grammar, productivity,
discriminative power. That does not make it a physical theory. It
does mean that the substrate's claims can be assessed at the
\emph{linguistic} level --- which dictionary entries are
well-defined, which grammar rules are tight, which translations
clear their bridge hypotheses --- rather than at the
``recurrence-by-construction'' level. We invite this kind of
engagement.

The paper does not change a single claim of the underlying
programme. It changes the form in which the programme's claims can
be reviewed.

\bigskip
\hrule
\bigskip

\noindent
\textbf{A reading note.}
The dictionary, grammar, and translation tables in the appendices
are typeset summaries of canonical CSVs in \texttt{data/}.
Disagreement with a row of any table is welcome; the CSVs are the
artefacts to file an issue against.

% =====================================================================
\appendix

\section{Dictionary table}
\label{app:dictionary}

The canonical machine-readable dictionary is at
\texttt{data/dictionary.csv}. The columns are:

\begin{enumerate}[leftmargin=*]
  \item \texttt{entry\_id} --- short identifier.
  \item \texttt{substrate\_object} --- the substrate primitive.
  \item \texttt{mathematical\_referent} --- the mathematical
        referent of the entry.
  \item \texttt{candidate\_physical\_referent} --- the candidate
        physical referent (if any).
  \item \texttt{candidate\_cognitive\_referent} --- the candidate
        cognitive referent (if any).
  \item \texttt{layer} --- 1, 2, or 3.
  \item \texttt{evidential\_status} --- \texttt{theorem-grade},
        \texttt{conditional}, \texttt{interpretive}.
  \item \texttt{bridge\_hypothesis} --- the named bridge required
        to lift the entry to Layer 3.
  \item \texttt{evidence\_ledger\_row} --- the
        \cite{evidenceLedger} row that audits this entry.
\end{enumerate}

The dictionary contains 14 entries in v1.0.0-rc1.

\section{Grammar table}
\label{app:grammar}

The canonical machine-readable grammar is at
\texttt{data/grammar.csv}. Columns:

\begin{enumerate}[leftmargin=*]
  \item \texttt{rule\_id}
  \item \texttt{transformation\_kind} --- DUAL,
        SHELL\_PROJECTION, ..., DIMENSIONAL\_PROJECTION (12 kinds).
  \item \texttt{input\_typing}
  \item \texttt{output\_typing}
  \item \texttt{preserved\_invariants}
  \item \texttt{broken\_invariants}
  \item \texttt{exchanged\_invariants}
  \item \texttt{default\_admissibility\_class} --- \texttt{EXACT},
        \texttt{CANDIDATE}, \texttt{APPROXIMATE}, \texttt{BROKEN},
        \texttt{DEGENERATE}.
  \item \texttt{engine\_module} --- the engine module that
        implements the rule.
\end{enumerate}

The grammar contains 12 rules in v1.0.0-rc1, one per
transformation kind, plus three derived rules for the Schl\"afli
decomposition, the $\lambda = 12$ lift, and the $\lambda = 12$
uniqueness scan.

\section{Translation table}
\label{app:translation}

The canonical machine-readable translation table is at
\texttt{data/translation.csv}. Columns:

\begin{enumerate}[leftmargin=*]
  \item \texttt{translation\_id}
  \item \texttt{dictionary\_entry\_id}
  \item \texttt{grammar\_verdict} --- \texttt{EXACT},
        \texttt{CANDIDATE}, \texttt{APPROXIMATE}, etc.
  \item \texttt{bridge\_hypothesis} --- the named bridge.
  \item \texttt{interpretive\_claim} --- the candidate
        interpretive claim.
  \item \texttt{audit\_caveat} --- the conservative caveat
        attached by \texttt{physical\_interpretation.py} (or
        equivalent).
  \item \texttt{open\_problems} --- what would have to happen for
        the translation to upgrade.
\end{enumerate}

The translation table contains 14 entries in v1.0.0-rc1.

\section{Engine reproducibility}
\label{app:engine}

Every claim of grammar in \S\ref{sec:grammar} is reproducible from
modules in the closure transformation engine
\cite{transformationEngine}. Specifically:

\begin{description}[leftmargin=*]
  \item[Theorem \ref{thm:schlafli}.]
        \texttt{closure\_transform\_engine.decomposition} ---
        functions \texttt{certify\_v600\_construction},
        \texttt{certify\_v24\_subgroup},
        \texttt{certify\_right\_cosets},
        \texttt{certify\_each\_coset\_is\_24cell},
        \texttt{certify\_a5\_coset\_action}.
  \item[Theorem \ref{thm:lambda12}.]
        \texttt{closure\_transform\_engine.keystone} ---
        functions \texttt{coset\_partition},
        \texttt{build\_d1\_graphs},
        \texttt{lambda12\_lift},
        \texttt{exact\_certify\_lambda12\_lift},
        \texttt{exact\_a5\_characters}.
  \item[$\lambda = 12$ uniqueness scan.]
        \texttt{closure\_transform\_engine.physics.universality\_lambda12}.
  \item[Transformation taxonomy and admissibility classes.]
        \texttt{closure\_transform\_engine.transformations.transformation\_grammar},
        \texttt{closure\_transform\_engine.physics.admissibility}.
  \item[Transition rules.]
        \texttt{closure\_transform\_engine.physics.transition\_rules}.
  \item[Physical-interpretation labels.]
        \texttt{closure\_transform\_engine.physics.physical\_interpretation}.
  \item[Compositional reasoning.] \texttt{closure\_ai}
        (relational architecture); \texttt{min\_aria}
        (vector-state agent).
\end{description}

All modules are pure Python with \texttt{networkx}, \texttt{numpy},
\texttt{scipy} dependencies. Exact $\Q(\sqrt 5)$ arithmetic is
provided by the vendored \texttt{vfd\_v600} package.

\bigskip

\bibliographystyle{plain}
\bibliography{references}

\end{document}
